\begin{theorem}
For any $s\geq3$, there is a positive constant $c_s$ such that for every
$k\geq2$,
\[
 r(s,k)\geq c_s\frac{k^{s-1}}{(\log k)^{2s-4}}.
\]
\end{theorem}
\begin{proof}
Fix $s\geq3$.  We first treat sufficiently large $k$.  Put $t=s-1$, and
let $C=C(t)$ be supplied by \noderef{DStarCounting}.  Choose the power of two
$q$ by the following exact dyadic rule.  Since $C$ is a positive integer,
there is a largest integer $m\geq0$ such that
\[
             C2^m m^2\leq k.                                  \tag{1}
\]
Indeed, the left side is eventually strictly increasing and tends to
infinity.  For all sufficiently large $k$, this maximal $m$ is at least
both $6$ and an integer $m_C$ with $2^{m_C}\geq C$.  Put $q=2^m$.
Then $q\geq C$, $q$ is a power of two, and
\[
\operatorname{Nat.log}_2 q=m,\qquad
 Cq(\operatorname{Nat.log}_2 q)^2=C2^m m^2\leq k.                \tag{2}
\]
Thus the exact hypotheses of \noderef{DStarCounting} hold.  Maximality in
(1) also gives
\[
 k<C2^{m+1}(m+1)^2=2Cq(m+1)^2.                                \tag{3}
\]
For $m\geq1$, (1) and $C\geq1$ imply $2^m\leq k$, hence
$m\leq\log_2 k$.  Therefore $m+1\leq2\log_2 k$, and (3) yields, with
$\log$ denoting the natural logarithm,
\[
 q>\frac{k}{2C(m+1)^2}
   \geq\frac{(\log 2)^2}{8C}\frac{k}{(\log k)^2}.             \tag{4}
\]
Here we used $\log_2 k=(\log k)/(\log2)$.  In particular, setting
$c'_s=(\log2)^2/(8C)>0$ gives the required explicit lower bound on $q$.

For $m\geq6$ the elementary inequality $(m+1)^2\leq2^m=q$ holds.
Combining this with (3), $t\geq2$, and $e>2$, we also have
\[
 k<2Cq(m+1)^2\leq2Cq^2\leq eCq^t.                            \tag{5}
\]
The cited lemma now supplies a loopless $T_s$-free digraph $D$ with
$|V(D)|\geq q^{2t-1}/4$ and
$\operatorname{fwi}_k(D)\leq(Cq^t)^k$.

Apply \noderef{RandomPermutationReduction} to obtain a loopless $K_s$-free
graph $G$ on the same vertex set.  Thus
\[
 i_k(G)\leq(e/k)^k(Cq^t)^k.
\]
Set
\[
                    p=\frac{k}{eCq^t}.
\]
All factors in this definition are positive, so (5) gives explicitly
$p>0$ and $p<1$, in particular $0\leq p\leq1$.  Multiplying the displayed
independent-set estimate by the nonnegative number $p^k$ gives
\[
 p^k i_k(G)
 \leq\left(\frac{k}{eCq^t}\right)^k
       \left(\frac{eCq^t}{k}\right)^k=1.                       \tag{6}
\]
Also $k\geq2$ gives $1\leq k$, while the construction gives a loopless
$K_s$-free graph and preserves the order of $D$.  Hence all hypotheses of
\noderef{SamplingDeletion} are now explicitly available.  That lemma yields
\[
 r(s,k)>p|V(G)|-1
 \geq\frac{k}{eCq^t}\frac{q^{2t-1}}4-1
 =\frac{kq^{t-1}}{4eC}-1.                                    \tag{7}
\]
Substituting (4) into (7), and using $t=s-1$, gives
\[
 r(s,k)>B_s\frac{k^{s-1}}{(\log k)^{2s-4}}-1,\qquad
 B_s=\frac{(c'_s)^{s-2}}{4eC}>0.                              \tag{8}
\]
For fixed $s\geq3$, the quantity
$X_k=k^{s-1}/(\log k)^{2s-4}$ tends to infinity.  We may therefore enlarge
the preceding threshold for $k$ so that $B_sX_k\geq2$.  Then the final
$-1$ is absorbed by the explicit inequality
\[
 B_sX_k-1\geq\frac{B_s}{2}X_k.                                \tag{9}
\]
Thus the desired lower bound holds for every sufficiently large $k$ with
the positive constant $B_s/2$.

There are only finitely many integers $k\geq2$ below this threshold.  By
\noderef{FiniteRamseyPositivity}, all corresponding $r(s,k)$ are positive.
For each of these finitely many $k$, define
\[
 \rho_{s,k}=r(s,k)\frac{(\log k)^{2s-4}}{k^{s-1}}.
\]
The positivity lemma and $k\geq2$ imply $\rho_{s,k}>0$.  Choose
\[
 c_s=\min\left(\frac{B_s}{2},\min_{2\leq k<K_s}\rho_{s,k}\right)>0,\tag{10}
\]
where $K_s$ is the large-$k$ threshold above.  For $k\geq K_s$, (9) and
$c_s\leq B_s/2$ prove the target inequality; for $2\leq k<K_s$,
$c_s\leq\rho_{s,k}$ proves it by the definition of $\rho_{s,k}$.
This establishes the displayed bound for every integer $k\geq2$.
\end{proof}
